% \begin{savequote}
% \begin{verbatim}
%  _________________________________________
% ( "Why make it simple when it is possible )
% ( to make it complicated"                 )
% (                                         )
% ( Husse March 2 2007                      )
%  -----------------------------------------
%        o   ,__,
%         o  (oo)____
%            (__)    )\
%               ||--|| *
% 
% \end{verbatim}
% \end{savequote}
\def\ChapterCode{MP1}
\chapter
[\CO{[MP1]} Medizinprodukte 1]
{Medizinprodukte 1}
\label{chp:GER}
\label{chp:MP1}

\ChapterIntro
{%
~
}
{%
\begin{description}
  \item[Maintainer:] Roman Koch
  \item[Autoren:] Michael Neuhold
  \item[Reviewer:] Standard-Reviewprozess
  \item[Version:] {\VarDocVersionTypeName} ({\VarDocVersionTypeDescription})
  \item[SHA1:] {\DocGitCommit}
\end{description}

{\TxTeilkapitelAASS}
}


\clearpage % TEMP temp clearpage


\section{Medizinproduktegesetz - MPG}
\label{topic:mpg-info}

\Layout{61}{}{%
\noindent Das Medizinproduktegesetz regelt die Funktionstüchtigkeit,
Sicherheit, Inbetriebnahme, Instandhaltung, den Betrieb und
die Anwendung, Überwachung, Sterilisation, Desinfektion und
Reinigung von Medizinprodukten. Als Medizinprodukte gelten
alle Instrumente, Apparate, Vorrichtungen, Stoffe oder
andere Gegenstände (inkl. Software), die vom Hersteller zur
Anwendung am Menschen zwecks Erkennen, Überwachen und Behandlung von
Krankheiten, Verletzungen oder Behinderungen bestimmt sind\ZFootnote{Genauer:
Zwecks
\begin{inparaenum}
  \item Erkennung, Verhütung, Überwachung, Behandlung oder Linderung von Krankheiten,
  \item Erkennung, Überwachung, Behandlung, Linderung oder Kompensierung von Verletzungen oder Behinderungen,
  \item Untersuchung, Veränderung oder zum Ersatz des anatomischen Aufbaus oder physiologischer Vorgänge oder
  \item Empfängnisregelung,
\end{inparaenum}
wobei \ldots{} Medikamente fallen nicht darunter.}.

Gegenstände und Stoffe, die dazu bestimmt sind, zusammen mit einem
Medizinprodukt verwendet zu werden (\enquote{Zubehör}), gelten selber auch als
Medizinprodukt. Dies trifft auch auf \E{Desinfektionsmittel} zu, mit welchen
Medizinprodukte gereinigt und wiederaufbereitet werden!


Es gibt tausende verschiedene Produkte, i.\,d.\,R. haben sie
alle eine \E{CE-Kennzeichnung}. Diese zeigt an, dass das
Produkt gewissen Minimalanforderungen entspricht. 
}
{}
{./images/geraete/Ce-logo.jpg}
{Das CE-Kennzeichen}
{---}




\Layout{57}{}
{}{}
{
\begin{tabulary}{\linewidth}{lLL}\toprule\addlinespace
%\addlinespace
\noindent\TabHeading{Klasse} &
\noindent\TabHeading{Beschreibung} &
\noindent\TabHeading{Beispiele}
\\\addlinespace\midrule\addlinespace
\TabSub{Klasse I}
&
Kein oder geringes Anwendungsrisiko
&
z.B: Verbandsmaterial, Gehhilfe, Rollstuhl, Tragsessel, Sehhilfe, Stethoskop, 
Trage, Beatmungsbeutel, Fieberthermometer, Druckminderer 
\\\addlinespace
\TabSub{Klasse IIa}
&
Anwendungsrisiko ohne entsprechender Ausbildung bzw. Unterweisung
&
z.B: aktive Diagnosegeräte, Absaugeinheit und -katheter, EKG-Kabel, 
Zuckermessstreifen, Trachealtuben, Dauerkatheter 
\\\addlinespace
\TabSub{Klasse IIb}
&
Erhöhtes Anwendungsrisiko ohne entsprechender Ausbildung bzw. Unterweisung
&
z.B: Beatmungsgerät, Defibrillator, Pulsoxymeter, Babyinkubator, Röntgengeräte
\\\addlinespace
\TabSub{Klasse III}
&
Besonders hohes Anwendungsrisiko ohne entsprechende Ausbildung bzw. Unterweisung
&
z.B: Perfusor, Insulinpumpe, Herzkatheter oder -klappen, Medizinprodukte mit Arzneimittelkomponenten
\\\addlinespace\bottomrule
\end{tabulary}
}
{MPG-Risikoklassen}{}

\TextSlide{Pflichten gemäß MPG}{
Anwender sind verpflichtet, die
Medizinprodukte richtig, zweckentsprechend und sicher
einzusetzen bzw.
anzuwenden, so wie es ihnen bei ihrer Einweisung gelehrt
wurde und es der Hersteller vorgesehen hat. Darüber hinaus
ergeben sich noch weitere spezielle Verpflichtungen:



\begin{itemize}
  \item \EE{Funktionskontrolle}:
  Der Anwender muss sich vor jeder Anwendung eines
  Medizinproduktes von der Funktionstüchtigkeit,
  Betriebssicherheit und dem ordnungsgemäßen Zustand überzeugen
  (\enquote{Gerätecheck}).
  \item \EE{Einweisung}:
  Medizinprodukte dürfen nur von solchen Personen angewendet
  werden, die auf Grund ihrer Ausbildung,  und
  erforderlichenfalls einer produktspezifischen \EE{Einweisung}
  das Medizinprodukt sachgerecht anwenden können und auch auf
  besondere anwendungs- und medizinproduktespezifische Gefahren
  hingewiesen worden sind. Dabei sind die Gebrauchsanweisungen
  sowie sonstige beigefügten Informationen der Produkte zu
  beachten und den Anwendern zugänglich zu machen.
  \item \E{Meldepflicht}:
  Fachpersonal muss Informationen über Medizinprodukte im
  Hinblick auf \E{Zwischenfälle}\footnote{Insbesonders jede
  Fehlfunktion, aber auch z.\,B. jeden Mangel in Bezug auf die
  Kennzeichnung oder die Gebrauchsanweisung, die zum Tod oder
  zu einer schwerwiegenden Schädigung eines Patienten,
  Anwenders oder eines Dritten führen können oder geführt
  haben}, oder schwerwiegende \E{Nebenwirkungen}, oder
  schwerwiegende \E{Qualitätsmängel}, unverzüglich dem
  Bundesamt für Sicherheit im Gesundheitswesen
  melden.\ZFootnote{\S\,70 MPG}
\end{itemize}
}
{
\begin{itemize}
  \item Sachgerechte Anwendung
  \item Einweisung
  \item Funktionskontrolle
  \item Meldepflicht bei schweren Mängeln
\end{itemize}
}
{}
{}
{}







%%%%%%%%%%%%%%%%%%%%%%%%%%%%%%%%%%%%%%%%%%%%%%%%%%%%%%%%%%%
%%%%%%%%%%%%%%%%%%%%%%%%%%%%%%%%%%%%%%%%%%%%%%%%%%%%%%%%%%%
%%%%%%%%%%%%%%%%%%%%%%%%%%%%%%%%%%%%%%%%%%%%%%%%%%%%%%%%%%%



\section{Steriles Material}

\Layout{22}{Entpacken von sterilem Material}
{
Die Verpackung von sterilem Material muss \E{fachkundig}
geöffnet werden, damit es beim Öffnen nicht verunreinigt
wird. Gängiges Verpackungsmaterial lässt sich
\EE{aufschälen}, d.\,h. die Oberseite lässt sich von der
Unterseite ablösen. Dabei darf das Material nicht
kontaminiert werden.

Manche vorgefertigten Sets (z.\,B. für diverse chirurgische
Eingriffe) werden nicht in foliierten Verpackungen
ausgeliefert, sondern sind in mehrere Lagen von Spezialpapier
eingewickelt. Das Spezialpapier darf nur an der Außenseite
angefasst werden und wird auf einer Fläche ausgebreitet,
sodass das darin eingewickelte Material steril bleibt.
}
{
\begin{itemize}
  \item Aufschälen
  \item Material muss steril bleiben!
\end{itemize}
}
{./images/AASdev20090228-img104.jpg}
{Aufschälen einer Verpackung}
{DEMAW}



% \begin{figure}[H]
% \includegraphics[width=6.656cm,height=4.993cm]{images/AASdev20090228-img103.jpg}
% \includegraphics[width=6.663cm,height=4.984cm]{images/AASdev20090228-img104.jpg}
% \Captionof{figure}{Öffnen von sterilen Material unter verschiedenen
% Bedingungen: Immer wird die Verpackung \enquote{aufgeschält}.
% \AuthNotesPicLegalNotOk{}{}}
% \end{figure}

% In diesem Abschnitt werden weitere Maßnahmen und Geräte für die Diagnostik und
% Patientenüberwachung (Monitoring) besprochen. Manche Maßnahmen werden in anderen
% Abschnitten besprochen (z.\,B. Blutdruckmessung
% (\prettyref{topic:blutdruckmessung}), Messung der Herzfrequenz (Pulsmessung,
% \prettyref{topic:pulsmessung}), \ldots). 

% \clearpage % TEMP temp clearpage

\section{Spezielle Medizinprodukte}

Für Medizinprodukte sind die Medizinprodukteeinweisungen
sowie die Gebrauchsanleitungen bzw. Herstellerhinweise zu
beachten. Allgemeine, nicht produktspezifische Informationen
finden sich in folgenden Abschnitten:

\begin{itemize}
  \item Absauggeräte \dotfill \FullPrettyRef{topic:absaugung}
  \item Beatmungshilfsmittel
  \begin{compactitem}
    \item Beatmungsbeutel \dotfill \FullPrettyRef{topic:beatmungsbeutel}
    \item Beatmungsgeräte \dotfill
    \FullPrettyRef{topic:beatmungsgeraete}
    \item PEEP-Ventil \dotfill \FullPrettyRef{topic:peep-ventil}
  \end{compactitem} 
  \item Berieselungshilfsmittel
  \begin{compactitem}
    \item Berieselungsbrille \dotfill \FullPrettyRef{topic:berieselungsbrille}
    \item Berieselungsmaske mit/ohne Reservoir \dotfill
    \FullPrettyRef{topic:berieselungsmaske}
  \end{compactitem}
  \item Sauerstoff \dotfill \FullPrettyRef{topic:sauerstoff}
  \item Geräte zum Retten von Patienten
  und Schienungsmaterial
  \begin{compactitem}
    \item Aluminiumkernschiene \dotfill
    \FullPrettyRef{topic:aluminiumkernschiene}
    \item HWS-Schiene \dotfill \FullPrettyRef{topic:hws-schiene-beschreibung}
    \item Rettungskorsett \dotfill \FullPrettyRef{topic:rettungskorsett-beschreibung}
    \item Schaufeltrage \dotfill \FullPrettyRef{topic:schaufeltrage-beschreibung}
    \item Spineboard \dotfill \FullPrettyRef{topic:spineboard-beschreibung}
    \item Vakuummatratze \dotfill \FullPrettyRef{topic:vakuummatratze-beschreibung}
    \item Vakuumschiene \dotfill \FullPrettyRef{topic:vakuumschiene}
  \end{compactitem}
  \item Transporthilfsmittel
  \begin{compactitem}
    \item (Fahr-)Trage \dotfill \FullPrettyRef{topic:fahrtrage}
    \item Tragering \dotfill \FullPrettyRef{topic:tragering}
    \item Tragetuch \dotfill \FullPrettyRef{topic:tragetuch-taetigkeit}
    \item Tragsessel \dotfill \FullPrettyRef{topic:tragsessel}
  \end{compactitem}
  \item Diagnostik und Monitoring
  \begin{compactitem}
    \item Blutdruckmessung \dotfill \FullPrettyRef{topic:blutdruckmessung}
    \item Blutzuckermessung \dotfill \FullPrettyRef{topic:blutzucker-messung}
    \item Pulsoxymetrie \dotfill \FullPrettyRef{topic:pulsoxymetrie}
    \item Stethoskop \dotfill \FullPrettyRef{topic:stethoskop}
    \item Temperaturmessung \dotfill \FullPrettyRef{topic:koerpertemperatur-messung}
  \end{compactitem}
\end{itemize}

%%%
\nocite{RueckerBildatlas1}
\ChapterEnd